\begin{textsample}{NEG Dim 5 – global\_north\_2024\_06 – Score -1.00 – global\_north\_2024\_06/109536945.txt}  \label{ex:f5_neg_089}
US military delivers supplies to Gaza via pier, but there ' s no one to distribute them U.S. officials are defending President Joe Biden ' s decision to build a \$230 million temporary pier off Gaza to deliver humanitarian aid as a better-than-nothing solution to the worsening crisis there, even as \textbf{rough} waves have forced its closure several times in the past month and aid workers on shore say it ' s nearly impossible to distribute the supplies due to safety concerns.
According to the military, which invited reporters to tour the pier for the first time on Tuesday, the structure has enabled the delivery of some 13.6 million pounds of food to a United Nations-run warehouse on shore.
The amount is roughly equivalent to what the U.N. says can fit inside 200 to 600 of its ground trucks -- less than half of what aid organizations say Gaza ' s 2 million residents might need in a day."This is a temporary solution to help rush aid into the zone, again recognizing the dire security situation there.
But we ' re going to continue to look at all @ @ @ @ @ @ @ @ @ @ Gen.
Pat Ryder, the Pentagon ' s press secretary.
Biden ' s decision to build the pier came amid frustration with the Israelis that they were allowing too few aid trucks inside Gaza, citing security concerns.
The plan was that the pier would enable some 2 million meals a day -- or 150 trucks -- that could augment ground convoys carrying supplies through border crossings like Kerem Shalom in Israel and Rafah in Egypt.
But that hasn't happened.
Since being anchored in mid-May, the pier has been operational about half the time it has been in place.
The temporary structure, built to rise and fall with the waves, has had to be moved twice to the Israeli port of Ashdod due to \textbf{rough} seas -- including once because portions of the pier broke and it had to be repaired.
In an interview with ABC News, Army Col.
Samuel Miller, commander of the 7th Transportation Brigade, said the hope is that seas are calmer in July.
He said it ' s possible though that @ @ @ @ @ @ @ @ @ @ into fall."We knew coming here there were going to be challenges.
And we ' ve seen just about every one of those challenges,"Miller told ABC News Chief Global Affairs Correspondent Martha Raddatz in an interview on Tuesday from the pier.
ABC News ' Martha Raddatz reports from the U.S. military pier floating off the coast of Gaza.
ABC News For now, he later added,"We ' re resolute and we ' re back out here."Officials say the military warned the White House beforehand of the likely complications, including high seas, that would make the pier inoperable.
But with few other options, Biden announced his decision during his State of the Union speech last March and ordered 1,000 U.S. troops to deploy to facilitate the maritime corridor.
A truck prepares to deliver aid to Gaza.
ABC News Aid organizations say an even bigger problem is that there ' s no system to distribute the aid once it arrives.
The U.N. ' s World Food @ @ @ @ @ @ @ @ @ @ the aid that arrives via the pier, as Biden has insisted that no U.S. troops deploy on shore in Gaza.
But WFP temporarily suspended operations earlier this month, pending a security review, and has not resumed operations.
Kate Phillips-Barrasso, vice president of global policy and advocacy for Mercy Corps, which operates inside Gaza, said the issue isn't with the pier itself but the"complete lack"of safety guarantees to aid workers from the Israelis."It doesn't matter if the aid comes in over land or through the sea.
It ' s impossible to deliver at a scale that would prevent massive food insecurity and potential famine"without guarantees given to aid workers and safe corridors, Phillips-Barrasso said.
The Gaza shoreline can be seen from the pier, with smoke visible.
ABC News When asked about the U.N. potentially suspending humanitarian operations across Gaza, the Pentagon ' s Ryder said there were no plans for U.S. service members to fill that role should that happen. @ @ @ @ @ @ @ @ @ @ organizations via USAID and other regional partners to ensure that we can find a way to get that,"he said.
In the meantime, Col.
Miller said the military troops tending to the pier will keep going until they are told to stop."We ' re out here to do our mission -- to move those pallets -- and those above us will make those decisions,"he said.

% matched lemmas: rough
\end{textsample}
